% Generated from locale/tr/content/incompleteness/representability-in-q/representable-comp.tex; do not hand-edit.


\olfileid[tr]{inc}{req}{rpc}
\olsection{$\Th{Q}$ İçinde Temsil Edilebilen Fonksiyonlar Hesaplanabilirdir}

$\Th{Q}$ içinde temsil edilebilen her fonksiyonun hesaplanabilir olduğunu
ispatlayacağız. Önce $\Th{Q}$ içinde temsil edilebilen fonksiyonlar hakkında
bir lemma kurmamız gerekir.

\begin{lem}\ollabel{lem:rep-q}
  $f(x_0, \dots, x_k)$, $\Th{Q}$ içinde temsil edilebiliyorsa öyle
  !!a{formula}~$!A(x_0, \dots, x_k, y)$ vardır ki
  \[
  \Th{Q} \Proves !A_f(\num{n_0}, \dots, \num{n_k}, \num{m})
  \quad\text{ancak ve ancak}\quad m = f(n_0, \dots, n_k).
  \]
\end{lem}

\begin{proof}
  ``Eğer'' yönü
\olref[int]{defn:representable-fn}\olref[int]{defn:rep:a}'dır. ``Yalnızca
eğer'' yönü şöyle görülür: $\Th{Q} \Proves !A_f(\num{n_0},
\dots, \num{n_k}, \num{m})$ fakat $m \neq f(n_0, \dots, n_k)$ olduğunu
varsayın. $l=f(n_0,\dots,n_k)$ olsun.
\olref[int]{defn:representable-fn}\olref[int]{defn:rep:a} uyarınca
$\Th{Q} \Proves !A_f(\num{n_0}, \dots, \num{n_k}, \num{l})$ olur.
\olref[int]{defn:representable-fn}\olref[int]{defn:rep:b} uyarınca
$\lforall[y][(!A_f(\num{n_0}, \dots, \num{n_k}, y) \lif \num{l} =
y)]$ olur. Mantığı ve $\Th{Q} \Proves
!A_f(\num{n_0}, \dots, \num{n_k}, \num{m})$ varsayımını kullanarak
$\Th{Q} \Proves \eq[\num{l}][\num{m}]$ elde ederiz. Öte yandan
\olref[bre]{lem:q-proves-neq} uyarınca $\Th{Q} \Proves
\eq/[\num{l}][\num{m}]$ olur. Böylece $\Th{Q}$ tutarsızdır. Fakat bu
olanaksızdır; çünkü $\Th{Q}$ standart modelde sağlanır (bkz.
\olref[int][def]{def:standard-model}), $\Sat{N}{\Th{Q}}$ olur ve
Sağlamlık Teoremi uyarınca sağlanabilir kuramlar her zaman tutarlıdır
(\tagrefs{prfAX/{fol:axd:sou:cor:consistency-soundness},
prfSC/{fol:seq:sou:cor:consistency-soundness},
prfND/{fol:ntd:sou:cor:consistency-soundness},
prfTab/{fol:tab:sou:cor:consistency-soundness}}).
\end{proof}

\begin{lem}
$\Th{Q}$ içinde temsil edilebilen her fonksiyon hesaplanabilirdir.
\end{lem}

\begin{proof}
Önce bunun neden doğru olduğuna ilişkin sezgisel düşünceyi verelim. $f$'yi
hesaplamak için şunları yaparız. Aritmetik dilindeki olası bütün
!!{derivation}s{}i~$\delta$ listeleyin. Bu mekanik olarak yapılabilir.
Her biri için bunun !!a{derivation} olup olmadığını denetleyin; türettiği
şey $!A_f(\num{n_0}, \dots, \num{n_k}, \num{m})$ biçimindeki
!!a{formula} olmalıdır (bu, \olref{lem:rep-q}'dan, $\Th{Q}$ içinde $f$'yi
temsil eden !!{formula}{}dür). Öyleyse \olref{lem:rep-q} uyarınca
$m=f(n_0,\dots,n_k)$'dır ve $f$'nin değerini bulmuşuzdur. Arama sona erer;
çünkü $\Th{Q} \Proves !A_f(\num{n_0}, \dots, \num{n_k},
\num{f(n_0, \dots, n_k)})$ olur ve dolayısıyla sonunda doğru türden bir
$\delta$ buluruz.

Yordamımız yalnız sayılar üzerinde değil, !!{derivation}s ve !!{formula}s
üzerinde çalıştığı ve ``olası bütün !!{derivation}s{}i listeleme''nin neden
mekanik olarak mümkün olduğunu tam olarak açıklamadığımız için bu henüz
kesin değildir. Fakat gördüğümüz üzere terimleri, !!{formula}s{}i ve
!!{derivation}s{}i Gödel sayılarıyla kodlamak mümkündür. Kesin bir hesaplama
modelini, genel özyinelemeli fonksiyonları da tanıttık. Ayrıca $d$,
$\Th{Q}$'nun aksiyomlarından !!a{derivation}{}in Gödel numarasıysa ve bu
türetimin sonucu Gödel numarası $y$ olan !!{formula} ise ancak ve ancak geçerli olan
$\Prf[\Th{Q}](d,y)$ bağıntısının (ilkel) özyinelemeli olduğunu gördük.
Gereksineceğimiz öteki ilkel özyinelemeli fonksiyonlar $\fn{num}$
(\olref[art][trm]{prop:num-primrec}) ve $\fn{Subst}$'tur
(\olref[art][sub]{prop:subst-primrec}). Bunlardan $f$'yi minimizasyonla
tanımlamak mümkündür; dolayısıyla $f$ özyinelemelidir.

Önce şunu tanımlayın:
\begin{multline*}
  A(n_0, \dots, n_k, m) = \\
  \fn{Subst}(\fn{Subst}(\dots\fn{Subst}(\Gn{!A_f}, \fn{num}(n_0), \Gn{x_0}),\\ \dots),
  \fn{num}(n_k),  \Gn{x_k}), \fn{num}(m), \Gn{y})
\end{multline*}
Bu karmaşık görünür; fakat yalnızca $A(n_0, \dots, n_k,
m) = \Gn{!A_f(\num{n_0}, \dots, \num{n_k}, \num{m})}$ fonksiyonudur.

Şimdi $(s)_0$, $\Th{Q}$'dan
$!A_f(\num{n_0}, \dots, \num{n_k}, \num{(s)_1})$'in !!a{derivation}{}inin
Gödel numarası olduğunda geçerli olan $R(n_0,\dots,n_k,s)$ bağıntısını
ele alın:
\[
R(n_0, \dots, n_k, s) \quad\text{ancak ve ancak}\quad \Prf[\Th{Q}]((s)_0, A(n_0,
\dots, n_k, (s)_1))
\]
$R(n_0,\dots,n_k,s)$'nin geçerli olduğu bir $s$ bulabilirsek,
$(s)_0$'ın $A_f(\num{n_0}, \dots, \num{n_k}, (s)_1)$'in
!!a{derivation}{}inin Gödel numarası olduğu bir sayı çifti, $(s)_0$ ve
$(s)_1$ bulmuş oluruz. Dolayısıyla $s$'yi aramak, biçimsel olmayan ispattaki
$d$ ile $m$ çiftini aramak gibidir. Böyle bir $s$'yi ``arayan'' hesaplanabilir
bir fonksiyon düzenli minimizasyonla tanımlanabilir. $R$'nin düzenli olduğuna
dikkat edin: her $n_0$, \dots, $n_k$ için
$\Th{Q} \Proves !A_f(\num{n_0}, \dots, \num{n_k},
\num{f(n_0, \dots, n_k)})$ ifadesinin !!a{derivation}{}i~$\delta$ vardır;
dolayısıyla $s=\tuple{\Gn{\delta},f(n_0,\dots,n_k)}$ için
$R(n_0,\dots,n_k,s)$ geçerlidir. Böylece $f$'yi şöyle yazabiliriz:
\[
f(n_0,\dots,n_{k}) = (\umin{s}{R(n_0, \dots, n_k, s)})_1.
\]
\end{proof}

